% GENERATED FILE. Edit canonical lessons, then run npm run generate:books.
% canonical-source-hash: fnv1a64:3892075670a3206c
% canonical-lessons: SA-C527-dhanurdharah, SA-C527-ekagrah, SA-C527-prahinoti, SA-C527-utsahayati, SA-C527-jya

\chapter{Aim and Shoot}
\label{ch:sa-aim-and-shoot}
\hlchaptermodality{527}

\begin{chapteropening}
\textbf{By the end of this chapter:} \emph{I can talk about archery, aim and encouragement.}

Complete the last lesson of chapter 527: I can talk about archery, aim and encouragement.
\end{chapteropening}

\section[\texorpdfstring{\sk{धनुर्धरः} (dhanurdharaḥ)}{धनुर्धरः (dhanurdharaḥ)}]{\sk{धनुर्धरः} (dhanurdharaḥ) — an archer}

\label{lesson:SA-C527-dhanurdharah}



\begin{quote}
Before the new one: say the Sanskrit for leaps over, then the Sanskrit for a team.
\end{quote}

\subsection*{You'll want to know: \sk{धनुर्धरः}}
\textbf{\sk{धनुर्धरः}} — \emph{dhanurdharaḥ} — \textquotedblleft{}an archer\textquotedblright{}.

\begin{quote}
\sk{धनुर्धरः} \sk{धनुः} \sk{गृह्णाति।}
\end{quote}

\begin{quote}
\emph{dhanurdharaḥ dhanuḥ gṛhṇāti.}
\end{quote}

\begin{quote}
The archer takes up his bow.
\end{quote}

A masculine noun.

\begin{cousinweb}
\textbf{\sk{धनुर्धरः}} (\emph{dhanurdharaḥ}) is built from \emph{dhanus}, \textquotedblleft{}a bow\textquotedblright{}, and \emph{dhara}, \textquotedblleft{}holding\textquotedblright{}, from the root √dhṛ, \textquotedblleft{}to hold\textquotedblright{}; MW gives \emph{dhanurdhara} as a bow-holder, an archer (MW).
\end{cousinweb}

\begin{grammarlens}[title={What you've built}]
One more word for describing what you lived and felt.
\end{grammarlens}

\subsection*{Your turn}
\begin{itemize}
\raggedright
  \item \emph{Say it:} \emph{dhanurdharaḥ}
  \item \emph{Say it:} \emph{dhanurdharaḥ}, once more
  \item \emph{Say it:} \emph{dhanurdharaḥ dhanuḥ gṛhṇāti.} — The archer takes up his bow.
  \item \emph{Recall:} say the Sanskrit for leaps over, then the Sanskrit for a team, then say \emph{dhanurdharaḥ} again
\end{itemize}

\begin{tcolorbox}[breakable,colback=teal!4,colframe=teal!35!black,title={Before you move on}]
What does \sk{धनुर्धरः} mean? (\textquotedblleft{}An archer\textquotedblright{}.) Say it once more.
\end{tcolorbox}
\section[\texorpdfstring{\sk{एकाग्रः} (ekāgraḥ)}{एकाग्रः (ekāgraḥ)}]{\sk{एकाग्रः} (ekāgraḥ) — focused; intent}

\label{lesson:SA-C527-ekagrah}



\begin{quote}
Before the new one: say the Sanskrit for a team, then the Sanskrit for an archer.
\end{quote}

\subsection*{You'll want to know: \sk{एकाग्रः}}
\textbf{\sk{एकाग्रः}} — \emph{ekāgraḥ} — \textquotedblleft{}focused; intent\textquotedblright{}.

\begin{quote}
\sk{धनुर्धरः} \sk{एकाग्रः} \sk{तिष्ठति।}
\end{quote}

\begin{quote}
\emph{dhanurdharaḥ ekāgraḥ tiṣṭhati.}
\end{quote}

\begin{quote}
The archer stands focused.
\end{quote}

\begin{cousinweb}
\textbf{\sk{एकाग्रः}} (\emph{ekāgraḥ}) is built from \emph{eka}, \textquotedblleft{}one\textquotedblright{}, and \emph{agra}, \textquotedblleft{}a point, a tip\textquotedblright{}; MW gives \emph{ekāgra} as fixing one's attention on one point, intent (MW).
\end{cousinweb}

\begin{grammarlens}[title={What you've built}]
One more word for describing what you lived and felt.
\end{grammarlens}

\subsection*{Your turn}
\begin{itemize}
\raggedright
  \item \emph{Say it:} \emph{ekāgraḥ}
  \item \emph{Say it:} \emph{ekāgraḥ}, once more
  \item \emph{Say it:} \emph{dhanurdharaḥ ekāgraḥ tiṣṭhati.} — The archer stands focused.
  \item \emph{Recall:} say the Sanskrit for a team, then the Sanskrit for an archer, then say \emph{ekāgraḥ} again
\end{itemize}

\begin{tcolorbox}[breakable,colback=teal!4,colframe=teal!35!black,title={Before you move on}]
What does \sk{एकाग्रः} mean? (\textquotedblleft{}Focused; intent\textquotedblright{}.) Say it once more.
\end{tcolorbox}
\section[\texorpdfstring{\sk{प्रहिणोति} (prahiṇoti)}{प्रहिणोति (prahiṇoti)}]{\sk{प्रहिणोति} (prahiṇoti) — shoots; sends forth}

\label{lesson:SA-C527-prahinoti}



\begin{quote}
Before the new one: say the Sanskrit for an archer, then the Sanskrit for focused.
\end{quote}

\subsection*{You'll want to know: \sk{प्रहिणोति}}
\textbf{\sk{प्रहिणोति}} — \emph{prahiṇoti} — \textquotedblleft{}shoots; sends forth\textquotedblright{}.

\begin{quote}
\sk{धनुर्धरः} \sk{लक्ष्यं} \sk{प्रति} \sk{बाणं} \sk{प्रहिणोति।}
\end{quote}

\begin{quote}
\emph{dhanurdharaḥ lakṣyaṃ prati bāṇaṃ prahiṇoti.}
\end{quote}

\begin{quote}
The archer shoots an arrow at the target.
\end{quote}

Past (he …): \sk{सः} \sk{बाणं} \sk{प्रहितवान्}, \emph{saḥ bāṇaṃ prahitavān}.

\begin{cousinweb}
\textbf{\sk{प्रहिणोति}} (\emph{prahiṇoti}) is built on \emph{pra}, \textquotedblleft{}forth\textquotedblright{}, and the root √hi, \textquotedblleft{}to send, to impel\textquotedblright{}; MW gives \emph{prahi} as to urge on, to send, to hurl; its past participle is \emph{prahita} (MW).
\end{cousinweb}

\begin{grammarlens}[title={What you've built}]
One more word for describing what you lived and felt.
\end{grammarlens}

\subsection*{Your turn}
\begin{itemize}
\raggedright
  \item \emph{Say it:} \emph{prahiṇoti}
  \item \emph{Say it:} \emph{prahiṇoti}, once more
  \item \emph{Say it:} \emph{dhanurdharaḥ lakṣyaṃ prati bāṇaṃ prahiṇoti.} — The archer shoots an arrow at the target.
  \item \emph{Recall:} say the Sanskrit for an archer, then the Sanskrit for focused, then say \emph{prahiṇoti} again
\end{itemize}

\begin{tcolorbox}[breakable,colback=teal!4,colframe=teal!35!black,title={Before you move on}]
What does \sk{प्रहिणोति} mean? (\textquotedblleft{}Shoots; sends forth\textquotedblright{}.) Say it once more.
\end{tcolorbox}
\section[\texorpdfstring{\sk{उत्साहयति} (utsāhayati)}{उत्साहयति (utsāhayati)}]{\sk{उत्साहयति} (utsāhayati) — cheers on; encourages}

\label{lesson:SA-C527-utsahayati}



\begin{quote}
Before the new one: say the Sanskrit for focused, then the Sanskrit for shoots.
\end{quote}

\subsection*{You'll want to know: \sk{उत्साहयति}}
\textbf{\sk{उत्साहयति}} — \emph{utsāhayati} — \textquotedblleft{}cheers on; encourages\textquotedblright{}.

\begin{quote}
\sk{गुरुः} \sk{शिष्यम्} \sk{उत्साहयति।}
\end{quote}

\begin{quote}
\emph{guruḥ śiṣyam utsāhayati.}
\end{quote}

\begin{quote}
The teacher encourages the pupil.
\end{quote}

Past (he …): \sk{सः} \sk{शिष्यम्} \sk{उत्साहितवान्}, \emph{saḥ śiṣyam utsāhitavān}.

\begin{cousinweb}
\textbf{\sk{उत्साहयति}} (\emph{utsāhayati}) is built on \emph{ud}, \textquotedblleft{}up\textquotedblright{}, and the causative of the root √sah, \textquotedblleft{}to bear, to be able\textquotedblright{}; MW gives \emph{utsāhayati} as to animate, to encourage (MW).
\end{cousinweb}

\begin{grammarlens}[title={What you've built}]
One more word for describing what you lived and felt.
\end{grammarlens}

\subsection*{Your turn}
\begin{itemize}
\raggedright
  \item \emph{Say it:} \emph{utsāhayati}
  \item \emph{Say it:} \emph{utsāhayati}, once more
  \item \emph{Say it:} \emph{guruḥ śiṣyam utsāhayati.} — The teacher encourages the pupil.
  \item \emph{Recall:} say the Sanskrit for focused, then the Sanskrit for shoots, then say \emph{utsāhayati} again
\end{itemize}

\begin{tcolorbox}[breakable,colback=teal!4,colframe=teal!35!black,title={Before you move on}]
What does \sk{उत्साहयति} mean? (\textquotedblleft{}Cheers on; encourages\textquotedblright{}.) Say it once more.
\end{tcolorbox}
\section[\texorpdfstring{\sk{ज्या} (jyā)}{ज्या (jyā)}]{\sk{ज्या} (jyā) — a bowstring}

\label{lesson:SA-C527-jya}



\begin{quote}
Before the new one: say the Sanskrit for shoots, then the Sanskrit for cheers on.
\end{quote}

\subsection*{You'll want to know: \sk{ज्या}}
\textbf{\sk{ज्या}} — \emph{jyā} — \textquotedblleft{}a bowstring\textquotedblright{}.

\begin{quote}
\sk{धनुर्धरः} \sk{ज्यां} \sk{कर्षति।}
\end{quote}

\begin{quote}
\emph{dhanurdharaḥ jyāṃ karṣati.}
\end{quote}

\begin{quote}
The archer draws the bowstring.
\end{quote}

A feminine noun.

\begin{cousinweb}
\textbf{\sk{ज्या}} (\emph{jyā}) is the word MW records for a bowstring, and MW compares Greek \emph{bios}, \textquotedblleft{}a bow\textquotedblright{} (MW).
\end{cousinweb}

\begin{grammarlens}[title={What you've built}]
One more word for describing what you lived and felt.
\end{grammarlens}

\subsection*{Your turn}
\begin{itemize}
\raggedright
  \item \emph{Say it:} \emph{jyā}
  \item \emph{Say it:} \emph{jyā}, once more
  \item \emph{Say it:} \emph{dhanurdharaḥ jyāṃ karṣati.} — The archer draws the bowstring.
  \item \emph{Recall:} say the Sanskrit for shoots, then the Sanskrit for cheers on, then say \emph{jyā} again
\end{itemize}

\begin{tcolorbox}[breakable,colback=teal!4,colframe=teal!35!black,title={Before you move on}]
What does \sk{ज्या} mean? (\textquotedblleft{}A bowstring\textquotedblright{}.) Say it once more.
\end{tcolorbox}
